% BHU PYQ -- Professional Exam Template
\documentclass[11pt,a4paper]{article}

\usepackage[a4paper,top=2.5cm,bottom=2.5cm,left=2.8cm,right=2.8cm,headheight=14pt]{geometry}
\usepackage{amsmath,amssymb}
\usepackage[T1]{fontenc}
\usepackage[utf8]{inputenc}
\usepackage{lmodern}
\usepackage{microtype}
\usepackage{fancyhdr}
\usepackage{enumitem}
\usepackage{hyperref}
\usepackage{lastpage}
\usepackage{parskip}

\hypersetup{colorlinks=true,urlcolor=black,linkcolor=black,
  pdftitle={B.Sc. (Hons.) Botany Sem-V BOB-502 -- BHU PYQ}}

\pagestyle{fancy}
\fancyhf{}
\renewcommand{\headrulewidth}{0.4pt}
\renewcommand{\footrulewidth}{0.4pt}
\fancyhead[L]{\small   \textit{Banaras Hindu University}}
\fancyhead[R]{\small   \textit{Botany Paper: BOB-502}}
\fancyfoot[C]{\small Page  \thepage\ of \pageref{LastPage}}


\newlist{parts}{enumerate}{2}
\setlist[parts,1]{label=(\alph*),leftmargin=*,itemsep=4pt,topsep=3pt}
\setlist[parts,2]{label=(\roman*),leftmargin=*,itemsep=2pt,topsep=2pt}

\newcommand{\pts}[1]{\hfill   \textit{[#1]}}


\begin{document}


\begin{center}
  {\large\bfseries BANARAS HINDU UNIVERSITY}\\[2pt]
  {
\normalsize B.Sc. (Hons.) Semester V Examination, 2013-14}\\[6pt]
  \rule{0.8\linewidth}{0.5pt}\\[6pt]
  {\large\bfseries Botany}\\[4pt]
  {\large\bfseries Paper: BOB-502 --- Comparative Studies of Phanerogams}\\[4pt]
  {
\normalsize Time: Three hours \hfill Max. Marks: 70}
\end{center}

\rule{\linewidth}{0.4pt}

\smallskip



\noindent   \textit{Note: Attempt any five questions selecting two questions from each Section A and B. Question No. 9 (Section-C) is compulsory. The figures in the right-hand margin indicate marks.}

\bigskip

\begin{center}
   \textbf{SECTION -- A}
\end{center}

\medskip


\noindent   \textbf{Question 1.} Give a detailed account of the order Gnetales with special reference to   \textit{Gnetum}.\pts{15}

\medskip


\noindent   \textbf{Question 2.} Describe briefly any two of the following:\pts{$7.5  \times 2 = 15$}

\begin{parts}
  \item Salient features of Ginkgoales
  \item   \textit{Williamsonia}
  \item Cycadales
\end{parts}

\medskip


\noindent   \textbf{Question 3.} Give a brief account of the economic importance of gymnosperms.\pts{15}

\medskip


\noindent   \textbf{Question 4.} Write short notes on any two of the following:\pts{$7.5  \times 2 = 15$}

\begin{parts}
  \item   \textit{Pentoxylon}
  \item Sexual reproduction in   \textit{Biota}
  \item Embryogeny in conifers
\end{parts}

\bigskip

\begin{center}
   \textbf{SECTION -- B}
\end{center}

\medskip


\noindent   \textbf{Question 5.} With the help of floral formula and floral diagram describe the family Rutaceae. Give names of three economically important plants of the family.\pts{15}

\medskip


\noindent   \textbf{Question 6.} Write short notes on any two of the following:\pts{$7.5  \times 2 = 15$}

\begin{parts}
  \item Outline of Takhtajan System of classification
  \item Apomixis
  \item Chemotaxonomy
\end{parts}

\medskip


\noindent   \textbf{Question 7.} Write an account on embryo culture. Give its practical application.\pts{15}

\medskip


\noindent   \textbf{Question 8.} Differentiate between any three of the following pairs of families:\pts{$5  \times 3 = 15$}

\begin{parts}
  \item Asteraceae and Meliaceae
  \item Euphorbiaceae and Moraceae
  \item Liliaceae and Polygonaceae
  \item Convolvulaceae and Scrophulariaceae
\end{parts}

\bigskip

\begin{center}
   \textbf{SECTION -- C}
\end{center}

\medskip


\noindent   \textbf{Question 9.} Answer the following in not more than fifty words each:\pts{$1  \times 10 = 10$}

\begin{parts}
  \item Name an angiosperm-like character of   \textit{Gnetum}.
  \item Define Pycnoxylic wood.
  \item Define hard wood.
  \item What is a bract scale?
  \item Name a gymnosperm which is known as a living fossil.
  \item Name a plant in which ochreate stipule is found.
  \item Which type of placentation is found in Asteraceae?
  \item Define adventive embryony.
  \item Give the botanical name of a parasitic angiosperm.
  \item Give the botanical name of a medicinal plant of Convolvulaceae.
\end{parts}

\vfill
\rule{\linewidth}{0.4pt}

\begin{center}\small
\href{https://bhu-exam.vercel.app/}{Click Here}\\[4pt]\href{https://me-aryan.vercel.app/}{Click Here}\\[4pt]\href{https://quashan-me.vercel.app/}{Click Here}
\end{center}

\end{document}
